\exercisegroup

\begin{enumerate}
\def\labelenumi{\arabic{enumi})}
\item
  设 E 是复赋范空间，f 是其底实向量空间 \(E_{0}\) 上的一个对称双线性型，使得 \(f(x, x) = \|x\|^{2}\) 对所有 \(x \in E\) 成立。证明：存在且仅存在一个 E 上的埃尔米特半双线性型 g，使得 \(f(x, y) = \mathcal{R}g(x, y)\)（利用 V, p.~2 的公式 (5) 证明 \(f(x, iy) = -f(ix, y)\)），从而 \(g(x, x) = \|x\|^{2}\)。
\item
  设 E 是实或复赋范空间。假设对 E 中每个（R 上的）2 维向量子空间 P，都存在一个对称双线性型 \(f_{\mathbf{P}}\)（定义在 \(\mathbf{P} \times \mathbf{P}\) 上），使得 \(f_{\mathbf{P}}(x,x) = \| x\|^2\) 对所有 \(x \in \mathbf{P}\) 成立。证明 \(f_{\mathbf{P}}\) 的定义无歧义，且存在一个埃尔米特半双线性型 \(g\)（在 \(\mathbf{E} \times \mathbf{E}\) 上），使得对每个实平面 \(\mathbf{P} \subset \mathbf{E}\)，有 \(f_{\mathbf{P}}(x,y) = \mathcal{R}g(x,y)\)，从而 \(g(x,x) = \| x\|^2\)。（若 E 是实向量空间，注意对 E 的每一对点有 \(\| x - y\|^2 + \| x + y\|^2 = 2(\| x\|^2 + \| y\|^2)\)，并由此导出恒等式
\end{enumerate}

\[
\left\| x + y + z \right\| ^ {2} - \left\| x + y \right\| ^ {2} - \left\| y + z \right\| ^ {2} - \left\| z + x \right\| ^ {2} + \left\| x \right\| ^ {2} + \left\| y \right\| ^ {2} + \left\| z \right\| ^ {2} = 0;
\]

若 E 是复向量空间，则应用习题 1。）

\begin{enumerate}
\def\labelenumi{\arabic{enumi})}
\setcounter{enumi}{2}
\tightlist
\item
  设 \(\mathbf{E}\) 是一个实有限维向量空间，维数为 \(n, f\) 是 \(\mathbf{E}\) 上的一个正且非退化的对称双线性型，而 \(\mathbf{B}_f\) 是由关系式 \(f(x, x) \leqslant 1\) 所定义的有界凸集。若 \(\mathbf{a} = (a_1, \ldots, a_n)\) 是 \(\mathbf{E}\) 的一个基，\(\Delta\) 是 \(f\) 关于该基的判别式，则我们把 \(\mathbf{B}_f\) 关于 \(\mathbf{a}\) 的体积称为数 \(v_{\mathbf{a}}(f) = \gamma_n |\Delta|^{-1/2}\)，其中 \(\gamma_n = \pi^{n/2} / \Gamma\left(\frac{n}{2} + 1\right)\)。若 \(\mathbf{b} = (b_1, \ldots, b_n)\) 是 \(\mathbf{E}\) 的第二个基，且
\end{enumerate}

\[
a _ {1} \wedge a _ {2} \wedge \dots \wedge a _ {n} = \delta b _ {1} \wedge b _ {2} \wedge \dots \wedge b _ {n},
\]

我们有 \(v_{\mathbf{b}}(f) = |\delta| v_{\mathbf{a}}(f)\)。

\begin{enumerate}
\def\labelenumi{\alph{enumi})}
\item
  证明：若 \(f\) 与 \(g\) 是两个正的非退化的对称双线性型，满足 \(\mathbf{B}_f \subset \mathbf{B}_g\)（它等价于 \(g \leqslant f\)），则 \(v_{\mathbf{a}}(f) \leqslant v_{\mathbf{a}}(g)\)（考虑 \(E\) 的一个同时对 \(f\) 和 \(g\) 都正交的基）。
\item
  设 A 是 E 中的对称紧凸集，以 0 为内点。证明：在所有正的非退化对称双线性型 \(f\)（在 \(\mathbf{E}\) 上且满足 \(\mathbf{A} \subset \mathbf{B}_f\)）之中，存在且仅存在一个，使 \(\mathbf{B}_f\) 的体积（关于 \(\mathbf{E}\) 的给定基）尽可能地小。（为证唯一性，注意若 \(\mathbf{A}\) 含于 \(\mathbf{B}_f\) 和 \(\mathbf{B}_g\)，则它含于 \(\mathbf{B}_{(f + g)/2}\)，且我们有 \(v_{\mathbf{a}}((f + g)/2) \leqslant \frac{1}{2}(v_{\mathbf{a}}(f) + v_{\mathbf{a}}(g))\) 对每个基 \(\mathbf{a}\)（\(\mathbf{E}\) 的、同时对 \(f\) 与 \(g\) 正交的基）成立。）
\item
  设 A 是 E 中的对称紧凸集，以 0 为内点，并设 \(f\) 是满足 \(\mathbf{A} \subset \mathbf{B}_f\) 且使 \(\mathbf{B}_f\) 关于 E 的给定基有最小体积的正的非退化对称双线性型。证明：存在 E 中的点 \(x_1, \ldots, x_n, u_1, \ldots, u_n\)，它们具有下列性质：
\end{enumerate}

α) 对每个 k，我们有 \(x_{k} \in A\) 与 \(f(x_{k}, x_{k}) = 1\)。

β) 序列 \((u_{1},\dots,u_{n})\) 是 E 关于 \(f\) 的标准正交基。

\(\gamma)\) 若置 \(x_{k} = \sum_{j=1}^{n} a_{kj} u_{j}\)（其中 \(1 \leqslant k \leqslant n\)），则我们有 \(a_{kj} = 0\)（对 \(k < j\)）以及 \(a_{kk}^{2} \geqslant (n - k + 1)/n\)。（对 \(k\) 用归纳法论证。设 \(x_{1}, \ldots, x_{k}, u_{1}, \ldots, u_{k}\) 均已构造，设 \(P_{k}\) 是（关于 \(f\)）从 E 到由 \(u_{1}, \ldots, u_{k}\) 生成的子空间上的正交投影算子；对每个 \(\varepsilon > 0\)，考虑由下式定义的双线性型 \(f_{\varepsilon}\)

\[
f _ {\varepsilon} (x, y) = (1 + \varepsilon) ^ {k - n} f \left(P _ {k} x, P _ {k} y\right) + (1 + \varepsilon) ^ {k} f \left(x - P _ {k} x, y - P _ {k} y\right)
\]

并证明 \(\mathbf{A} \not\subset \mathbf{B}_{f_{\varepsilon}}\)（利用 \(b\) )）。对每个整数 \(p \geqslant 1\)，选取一点 \(y_{p}\)（在 \(\mathbf{A}\) 中、不属于 \(\mathbf{B}_{f_{1/p}}\)）；取 \(x_{k+1}\) 为序列 \((y_{p})\) 的一个极限点，使得

\[
k _ {f} \left(x _ {k + 1} - P _ {k} x _ {k + 1}, x _ {k + 1} - P _ {k} x _ {k + 1}\right) \geqslant (n - k) f \left(P _ {k} x _ {k + 1}, P _ {k} x _ {k + 1}\right);
\]

然后选取 \(u_{k+1}\)。

\(d\) ) 证明 \(b\) ) 与 \(c\) ) 关于满足 \(\mathbf{B}_f\subset \mathbf{A}\) 且使 \(\mathbf{B}_f\) 的体积（关于 E 的给定基）尽可能大的正的非退化对称双线性型的类似命题。

¶ 4) a) 设 E 是维数 ≥ 2 的实或复赋范空间，具有下列性质：关系式 \textbar\textbar x\textbar\textbar{} = \textbar\textbar y\textbar\textbar{} 蕴含不等式

\[
\left\| x + y \right\| ^ {2} + \left\| x - y \right\| ^ {2} \leqslant 2 \left(\left\| x \right\| ^ {2} + \left\| y \right\| ^ {2}\right).
\]

证明 E 上的范数是准希尔伯特的。（借助 V, p.~60 的习题 1 与习题 2，化归为 E 是实且 2 维的情形。此时设 A 是 E 的单位球，并设 f 是满足 \(A \subset B_{f}\) 且使 \(B_{f}\) 关于给定基的体积尽可能小的正的非退化对称双线性型。设 \(x_{1}, x_{2}\) 是习题 3, c) 中构造的两点。证明圆 \(f(z, z) = 1\) 与两向量 \(x_{1}, x_{2}\) 的分角线的交点也属于 A，并由迭代得出 \(A = B_{f}\)。）

\begin{enumerate}
\def\labelenumi{\alph{enumi})}
\setcounter{enumi}{1}
\item
  设 E 是维数 \(\geqslant 2\) 的实或复赋范空间，具有下列性质：关系 \(\| x + y \| = \| x - y \|\) 蕴含 \(\| x + y \|^2 = \| x \|^2 + \| y \|^2\)。证明 E 上的范数是准希尔伯特的（化归为 \(a\) )）。
\item
  证明 \(a\) 的类似命题，此时假设关系 \(\| x\| = \| y\|\) 蕴含不等式
\end{enumerate}

\[
\left\| x + y \right\| ^ {2} + \left\| x - y \right\| ^ {2} \geqslant 2 \left(\left\| x \right\| ^ {2} + \left\| y \right\| ^ {2}\right)
\]

（利用习题 3, d)）。

\begin{enumerate}
\def\labelenumi{\arabic{enumi})}
\setcounter{enumi}{4}
\tightlist
\item
  设 E 是维数 \(\geqslant 2\) 的实或复向量空间。假设给定了映射 \(x \mapsto \|x\|\)（从 E 到 \(R_{+}\) 中），它满足：\(\|\lambda x\| = |\lambda|\) · \(\|x\|\)（对每个标量 \(\lambda\)）、\(\|x\| = 0\) 蕴含 x = 0，以及下列“托勒密不等式”
\end{enumerate}

\[
\| a - c \| \cdot \| b - d \| \leqslant \| a - b \| \cdot \| c - d \| + \| b - c \| \cdot \| a - d \|
\]

对 E 中的所有 \(a, b, c, d\)。

\begin{enumerate}
\def\labelenumi{\alph{enumi})}
\item
  证明 \(\| x\|\) 是 E 上的范数（把 \(d\) 换成 0，把 \(b\) 换成 \(-a\)）。
\item
  证明 E 上的范数是准希尔伯特的。（由托勒密不等式推出不等式 \(\|x+y\|^{2}+\|x-y\|^{2}\geqslant4\|x\|. \|y\|\)，并利用习题 4, c)。）证明其逆（证明在希尔伯特空间中，若置 \(a' = a/\|a\|^{2}\)、\(b' = b/\|b\|^{2}\)，则有等式 \(\|a'-b'\|=\|a-b\|/\|a\|. \|b\|\)）。若 \(\|a\|=\|b\|=\|c\|=\|d\|\)，且四向量 a、b、c、d 位于同一平面内，则托勒密不等式的两端相等。
\end{enumerate}

\begin{enumerate}
\def\labelenumi{\arabic{enumi})}
\setcounter{enumi}{5}
\tightlist
\item
  \begin{enumerate}
  \def\labelenumii{\alph{enumii})}
  \tightlist
  \item
    设 E 是维数 \(\geqslant 2\) 的实或复赋范空间。证明：对 E 中每个 \(x \neq 0\) 及每个实数 \(\alpha > 0\)，存在 E 中的元素 \(y\)，使得 \(\| y \| = \alpha\) 且 \(\| x + y \|^2 = \| x \|^2 + \| y \|^2\)。
  \end{enumerate}
\end{enumerate}

\begin{enumerate}
\def\labelenumi{\alph{enumi})}
\setcounter{enumi}{1}
\item
  假设当向量 \(x, y\)（在 \(\mathbf{E}\) 中）满足关系 \(\| x + y\|^2 = \| x\|^2 + \| y\|^2\) 时，我们也有 \(\| x - y\|^2 = \| x\|^2 + \| y\|^2\)。证明 \(\mathbf{E}\) 上的范数是准希尔伯特的。（利用 a)，化归为习题 4 中的情形：限于 \(\mathbf{E}\) 是 2 维的情形，我们证明若 \(\| x_1\| = \| x_2\| = 1\)、\(y = \frac{1}{2}(x_1 - x_2)\)，且若 \(z \in \mathbf{E}\) 满足 \(\| y\|^2 + \| z\|^2 = \| y + z\|^2 = 1\)，则 \(z = \frac{1}{2}(x_1 + x_2)\) 或 \(z = -\frac{1}{2}(x_1 + x_2)\)。）
\item
  假设对每个向量 \(x \neq 0\)（在 \(E\) 中），集 \(H\)（由所有向量 \(y\) 中满足 \(\| x - y \|^2 = \| x \|^2 + \| y \|^2\) 者组成）在加法下稳定。证明 \(b\) ) 的结论成立。（化归为 \(E\) 是实且 2 维的情形。利用 \(a\) ) 以及 \(E\) 中单位球的紧性，证明 \(H\) 是一个闭集，它包含至少两条以 0 为原点的不同半直线；证明这两条半直线彼此相反，方法是注意：若不然，它们所生成的凸集将含于 \(H\)，并且将含有 \(x\) 或 \(-x\)。）
\end{enumerate}

\begin{enumerate}
\def\labelenumi{\arabic{enumi})}
\setcounter{enumi}{6}
\tightlist
\item
  设 E 是维数 \(\geqslant 2\) 的实或复赋范空间，具有下列性质：存在一个实数 \(\gamma\)（异于 0 与 \(\pm 1\)），使得关系 \(\| x + y \| = \| x - y \|\) 蕴含 \(\| x + \gamma y \| = \| x - \gamma y \|\)。
\end{enumerate}

\begin{enumerate}
\def\labelenumi{\alph{enumi})}
\item
  证明：若 \(\|x+y\|=\|x-y\|\)，则凸映射 \(\phi:\xi\mapsto\|x+\xi y\|\)（从 R 到 R）在任何区间上都不是常数。（用反证法论证；设 \([\alpha,\beta]\) 是使 \(\phi\) 为常数的最大区间，证明存在 \(\delta>\beta\)（足够接近 \(\beta\)）使 \(\phi(\delta)=\phi(\beta)\)，方法是注意：若 \(\|u+v\|=\|u-v\|\)，则对每个有理整数 n 有 \(\|u+\gamma^{n}v\|=\|u-\gamma^{n}v\|\)）。
\item
  证明：若 \(\|x + y\| = \|x - y\|\)，则有 \(\|x + \xi y\| = \|x - \xi y\|\)（对每个实数 \(\xi\)）。（用 a) 的记号，注意 \(\phi\) 在点 \(\xi = 0\) 处有一个相对极小，这利用了如下事实：对每个有理整数 n 有 \(\phi(\gamma^{n}) = \phi(-\gamma^{n})\)；由此推出我们恒有 \(\phi(\xi) = \phi(-\xi)\)，因为否则，会得到 \(\phi(\lambda) = \phi(\mu)\) 对两个数 \(\lambda, \mu\)（满足 \(\lambda + \mu \neq 0\)）成立，而此时 \(\phi\) 在点 \(\frac{1}{2}(\lambda + \mu)\) 处有一个相对极小。
\item
  由 b) 推出 E 上的范数是准希尔伯特的（首先证明：若 \(\|x\| = \|y\|\)，则我们有 \(\|\alpha x + \beta y\| = \|\beta x + 2y\|\)（对每对实数 \(\alpha\)、\(\beta\)），且等式 \(\|x + y\| = \|x - y\|\) 蕴含 \(\|\alpha x + \beta y\| = \|\alpha x - \beta y\|\)（对每对实数 \(\alpha\)、\(\beta\)）。其次证明：若 \(\|x\| = \|y\| = 1\) 且 \(\|x + y\| = \|x - y\|\)，则利用前面的结果可得关系式 \(\left\|\left(\alpha^{2} - \beta^{2}\right)x + 2\alpha\beta y\right\| = \alpha^{2} + \beta^{2}\)，并由此得出结论）。
\end{enumerate}

\begin{enumerate}
\def\labelenumi{\arabic{enumi})}
\setcounter{enumi}{7}
\tightlist
\item
  设 \(E\) 是实或复赋范空间，具有下列性质：若 \(x, y, x', y'\)（\(E\) 中的四个向量）满足
\end{enumerate}

\[
\| x \| = \| x ^ {\prime} \|, \quad \| y \| = \| y ^ {\prime} \|, \quad \| x + y \| = \| x ^ {\prime} + y ^ {\prime} \|,
\]

则 \(\| x - y \| = \| x' - y' \|\)。证明 \(E\) 上的范数是准希尔伯特的（利用习题 7）。

\begin{enumerate}
\def\labelenumi{\arabic{enumi})}
\setcounter{enumi}{8}
\item
  设 E 是实希尔伯特空间，f 是 E 上的连续线性型。证明：在 E 的每个闭凸子集 A 上，函数 \(x \mapsto \|x\|^{2} - f(x)\) 有下界，且在 A 的唯一一点处达到其极小。
\item
  设 \(\mathbf{E}\) 是实希尔伯特空间，\(\mathbf{B}\) 是 \(\mathrm{E} \times \mathrm{E}\) 上的双线性型，\(c_{1}\)、\(c_{2}\) 是两个数 \(> 0\)，使得
\end{enumerate}

\[
\left| \mathrm{B} (x, y) \right| \leqslant c _ {1} \| x \|. \| y \| \quad \text { for   every } x, y \text { in } E;
\]

\[
\left| \mathrm{B} (x, x) \right| \geqslant c _ {2} \| x \| ^ {2} \quad \text {   for   every   } x \in \mathrm{E}.
\]

证明：对每个连续线性型 \(f\)（在 \(\mathbf{E}\) 上），存在唯一的元素 \(x_{f} \in \mathbf{E}\)（相应地 \(y_{f} \in \mathbf{E}\)），使得 \(f(x) = \mathbf{B}(x_{f}, x)\)（相应地 \(f(x) = \mathbf{B}(x, y_{f})\)）对每个 \(x \in \mathbf{E}\) 成立。

\begin{enumerate}
\def\labelenumi{\arabic{enumi})}
\setcounter{enumi}{10}
\tightlist
\item
  设 \(\mathbf{E}\) 是希尔伯特空间，\((x_{n})\) 是 \(\mathbf{E}\) 中弱收敛于一点 \(a\) 的点序列。对每个 \(y \in \mathbf{E}\)，我们置
\end{enumerate}

\[
d (y) = \lim _ {n \rightarrow \infty} \inf \| x _ {n} - y \| \quad \text { and } \quad D (y) = \lim _ {n \rightarrow \infty} \sup \| x _ {n} - y \|.
\]

证明 \(d(y)^2 = d(a)^2 + \| y - a\|^2\) 与 \(\mathrm{D}(y)^2 = \mathrm{D}(a)^2 + \| y - a\|^2\)。若 \(\alpha\) 与 \(\beta\) 是满足 \(0 \leqslant \alpha \leqslant \beta\) 的两个实数，给出使 \(d(a) = \alpha\) 且 \(\mathrm{D}(a) = \beta\) 的例子。

¶ 12) a) 证明：存在一个数 \(c_{0} > 0\)，使得对每个维数为 \(n\) 的实赋范向量空间 E 以及每个整数 \(k \leqslant c_{0} n\)，存在 E 上的一个希尔伯特范数 \(x \mapsto \| x \|_{2}\) 使得 \(\| x \|_{2} \leqslant \| x \|\)（对所有 \(x \in \mathbf{E}\)），以及一个标准正交系 \((x_{j})_{1 \leqslant j \leqslant k}\)（由 E 的 \(k\) 个元素组成，关于该希尔伯特结构），其范数满足 \(\| x_{j} \| \leqslant 2\)。（利用 V, p.~60 的习题 3。）b) 设 \(n, m\) 是两个整数 \(> 0\)，满足 \(n \leqslant c_{0} m\)。设 E 是维数为 \(m\) 的实赋范向量空间。证明：存在 E 的一个维数为 \(n\) 的向量子空间 F、F 上的一个正且非退化的对称双线性型 \((x, y) \mapsto \langle x | y \rangle\)，以及 F 的一个标准正交基 \(\{a_{1}, a_{2}, ..., a_{n}\}\)，使得

\[
\frac {1}{2} \sup _ {j} | \langle a _ {j} | x \rangle | \leqslant \| x \| \leqslant \| x \| _ {2}
\]

（其中 \(\| x\| _2^2 = \langle x|x\rangle\)）对所有 \(x\in \mathrm{F}\)。（把 a) 应用于 E 的对偶 E'。）

\begin{enumerate}
\def\labelenumi{\arabic{enumi})}
\setcounter{enumi}{12}
\tightlist
\item
  \begin{enumerate}
  \def\labelenumii{\alph{enumii})}
  \tightlist
  \item
    设 \((x_{n})_{n\in \mathbb{N}}\) 是 Banach 空间 E 中的无穷序列。证明：为使族 \((x_{n})\) 可和，必须且只需：对每个序列 \((\varepsilon_{n})\)（由等于 1 或 \(-1\) 的数组成），以 \((\varepsilon_{n}x_{n})\) 为通项的级数收敛（利用 GT, III, § 5, 习题 4）。
  \end{enumerate}
\end{enumerate}

\begin{enumerate}
\def\labelenumi{\alph{enumi})}
\setcounter{enumi}{1}
\tightlist
\item
  设 \((x_{j})_{1\leqslant j\leqslant n}\) 是希尔伯特空间 \(E\) 中的有限点序列；证明
\end{enumerate}

\[
2 ^ {- n} \sum_ {(\varepsilon_ {j})} \left(\left\| \sum_ {j = 1} ^ {n} \varepsilon_ {j} x _ {j} \right\| ^ {2}\right) = \sum_ {j = 1} ^ {n} \| x _ {j} \| ^ {2},
\]

其中 \((\varepsilon_{j})\) 取遍由 \(2^{n}\) 个序列组成的集（这些序列的项都等于 1 或 \(-1\)）（利用中线恒等式，cf.~V, p.~9, 公式 (14)）。

\begin{enumerate}
\def\labelenumi{\alph{enumi})}
\setcounter{enumi}{2}
\tightlist
\item
  由 \(b)\) 推出：若 \((x_{i})_{i\in I}\) 是希尔伯特空间 \(E\) 中的可和族，则族 \((\| x_i\|^2)_{i\in I}\) 在 \(\mathbf{R}\) 中可和。
\end{enumerate}

¶ 14) 设 E 是无穷维 Banach 空间。a) 证明：对每个整数 N，存在由 E 中 N 个范数为 1 的向量组成的序列 \((b_{j})_{1 \leqslant j \leqslant N}\)，使得对每个由 N 个标量组成的序列 \((\xi_{j})_{1 \leqslant j \leqslant N}\)，我们有

\[
\left\| \sum_ {j = 1} ^ {\mathrm{N}} \xi_ {j} b _ {j} \right\| ^ {2} \leqslant 4 \sum_ {j = 1} ^ {\mathrm{N}} | \xi_ {j} | ^ {2}
\]

（利用习题 12, b)）。

\begin{enumerate}
\def\labelenumi{\alph{enumi})}
\setcounter{enumi}{1}
\item
  对每个数列 \((\lambda_n)_{n\geqslant 1}\)（由数 \(\geqslant 0\) 组成，满足 \(\sum_{n}\lambda_{n}^{2} < +\infty\)），证明：存在一个由 E 的点组成的序列 \((x_{n})_{n\geqslant 1}\)，使得 \(\| x_n\| = \lambda_n\)（对所有 \(n\)），且级数 \((x_{n})\) 可和。（利用习题 13, a) 的 a)。）
\item
  由 b) 推出：在每个无穷维 Banach 空间中，都存在一个交换收敛但不绝对收敛的级数（Dvoretzky-Rogers 定理）。
\end{enumerate}

\begin{enumerate}
\def\labelenumi{\arabic{enumi})}
\setcounter{enumi}{14}
\tightlist
\item
  设 E 是复希尔伯特空间，\(E_{1}\)、\(E_{2}\) 是 E 的两个闭向量子空间，\(P_{1}\)、\(P_{2}\) 分别是 E 到 \(E_{1}\)、\(E_{2}\) 上的正交投影算子。
\end{enumerate}

\begin{enumerate}
\def\labelenumi{\alph{enumi})}
\item
  证明：为使 \(P_{1}\) 与 \(P_{2}\) 交换，必须且只需 E 是四个子空间 \(E_{1} \cap E_{2}\)、\(E_{1}^{\circ} \cap E_{2}^{\circ}\)、\(E_{1}^{\circ} \cap E_{2}\)、\(E_{1} \cap E_{2}^{\circ}\) 的希尔伯特和（其中 \(M^{\circ}\) 表示 E 的向量子空间 M 的正交补）。
\item
  证明：若 \(E_{1}\) 有限维且 \(\|P_{1}-P_{2}\|<1\)，则 \(E_{2}\) 与 \(E_{1}\) 有相同的维数（考虑交 \(E_{1}^{\circ}\cap E_{2}\)）。
\item
  证明：E 的自同态 \(T=(P_{1}-P_{2})^{2}\) 与 \(P_{1}\) 和 \(P_{2}\) 交换，且 T 的对应于特征值 0（相应地 1）的特征子空间是正交子空间 \(E_{1}\cap E_{2}\) 与 \(E_{1}^{\circ}\cap E_{2}^{\circ}\)（相应地 \(E_{1}^{\circ}\cap E_{2}\) 与 \(E_{1}\cap E_{2}^{\circ}\)）的直和。
\item
  设 E 有限维且 \(T=\lambda I\)（其中 \(\lambda\neq0\)）。则 \(\lambda>0\)，且 E 是维数 \(\leqslant2\) 的子空间的希尔伯特和，其中每个子空间在 \(P_{1}\) 与 \(P_{2}\) 下稳定（注意 \(P_{1}-P_{2}\) 是自伴的，并由此推出 E 是两个子空间 \(E^{+},E^{-}\) 的希尔伯特和，使得 \(P_{1}.x-P_{2}.x=\sqrt{\lambda}x\)（在 \(E^{+}\) 中）且 \(P_{1}.x-P_{2}.x=-\sqrt{\lambda}x\)（在 \(E^{-}\) 中）；然后证明对 \(x\in E^{+}\)，有 \(P_{1}.x=\frac{1+\sqrt{\lambda}}{2}x+z\) 与 \(P_{2}.x=\frac{1-\sqrt{\lambda}}{2}x+z\)（其中 \(z\in E^{-}\)））。* e) 设 \(E_{1}\) 与 \(E_{2}\) 有限维。证明：存在 E 的一个子空间族 \((\mathrm{F}_{\alpha})_{\alpha\in\mathrm{A}}\)（子空间维数为 \(\leqslant2\)），使得 E、\(E_{1}\) 与 \(E_{2}\) 分别是族 \((\mathrm{F}_{\alpha})_{\alpha\in\mathrm{A}},(\mathrm{F}_{\alpha}\cap\mathrm{E}_{1})_{\alpha\in\mathrm{A}}\) 与 \((\mathrm{F}_{\alpha}\cap\mathrm{E}_{2})_{\alpha\in\mathrm{A}}\) 的希尔伯特和（利用 c) 化归为 E 有限维的情形，然后应用 d)）。*
\end{enumerate}

\begin{enumerate}
\def\labelenumi{\arabic{enumi})}
\setcounter{enumi}{15}
\tightlist
\item
  设 \(\mathbf{E}\) 是希尔伯特空间，\(P\) 是 \(\mathbf{E}\) 上的连续投影算子，即 \(\mathbf{E}\) 的满足 \(P^2 = P\) 的连续自同态。证明：为使 \(P\) 是正交投影算子，必须且只需 \(\| P \| \leqslant 1\)。（为看出该条件是充分的，考虑一个向量 \(x\)（与 \(I - P\) 的核正交）。）
\end{enumerate}

若 P 具有有限秩，证明：存在 E 的一个具有有限余维数的闭子空间 F，它包含 \(P(\mathrm{E})\)，且 P 在 F 上的限制是一个正交投影算子。

\begin{enumerate}
\def\labelenumi{\arabic{enumi})}
\setcounter{enumi}{16}
\tightlist
\item
  \begin{enumerate}
  \def\labelenumii{\alph{enumii})}
  \tightlist
  \item
    设 E 是 2 维实希尔伯特空间，\(P_1\)、\(P_2\) 分别是 E 到两条直线 \(\mathbf{D}_1\)、\(\mathbf{D}_2\)（假定不同）上的正交投影算子。证明：对每个 \(x \in E\)（满足 \(\| x \| = 1\)），我们有 \(\| (P_1 - P_2).x \| = \sin \theta\)，其中 \(\theta\) 是 \(\mathbf{D}_1\) 与 \(\mathbf{D}_2\) 之间介于 0 与 \(\pi/2\) 之间的夹角；且对 E 中每个 \(y \neq 0\)，存在 \(x \neq 0\) 使 \((P_1 - P_2).x\) 与 \(y\) 共线。
  \end{enumerate}
\end{enumerate}

\begin{enumerate}
\def\labelenumi{\alph{enumi})}
\setcounter{enumi}{1}
\item
  设 E 是实希尔伯特空间，\(P_{1}\)、\(P_{2}\) 是 E 上的两个正交投影算子，其像分别为 \(E_{1}\)、\(E_{2}\)。证明 \(\|P_{1}-P_{2}\|\) 是诸数 \(\sin\theta\) 的下确界，其中 \(\theta\) 是介于 0 与 \(\pi/2\) 之间的角，它是两条直线 \(D_{1}\)、\(D_{2}\) 的夹角（满足 \(D_{1}\subset E_{1}\)、\(D_{2}\subset E_{2}\)，且 \(D_{1}\) 与 \(D_{2}\) 均与 \(E_{1}\cap E_{2}\) 正交）。
\item
  设 \(Q_{1}, Q_{2}\) 是 \(\mathbf{E}\) 上的两个连续投影算子，其像为 \(\mathrm{E}_1, \mathrm{E}_2\)，并设 \(P_{1}, P_{2}\) 分别是到 \(\mathrm{E}_1\) 与 \(\mathrm{E}_2\) 上的正交投影算子。证明 \(\| P_2 - P_1 \| \leqslant \| Q_2 - Q_1 \|\)。（注意 \((Q_2 - Q_1)P_2 = (I - Q_1)(P_2 - P_1)\)，并利用 a) 与 b)。）
\end{enumerate}

¶ 18) 设 E 是维数 \(\geqslant 3\) 的实赋范空间。假设存在一个递减双射映射 \(\omega\)（从 E 的闭向量子空间所成之集 \(\mathfrak{M}\) 到其自身），使得 \(\omega(\omega(M)) = M\) 且 \(M \cap \omega(M) = \{0\}\)（对每个 \(M \in \mathfrak{M}\)）。

\begin{enumerate}
\def\labelenumi{\alph{enumi})}
\item
  证明：存在一个线性映射 \(u\)（从 \(\mathbf{E}\) 到其对偶 \(\mathbf{E}'\) 上），它在相差一个标量因子的意义下确定，且使得 \(u(\mathbf{M}) = (\omega(\mathbf{M}))^\circ\)（对所有 \(\mathbf{M} \in \mathfrak{M}\)）。（考虑 \(\mathbf{M}\) 是 1 维的情形，注意域 \(\mathbf{R}\) 的唯一自同构是恒等映射（GT, IV, § 3, 习题 3），应用射影几何的基本定理（A, II, § 9, 习题 16）。）
\item
  若置 \(\langle x|y\rangle = \langle x,u(y)\rangle\)，证明 \(\langle x|x\rangle \neq 0\)（对所有 \(x\neq 0\)），且关系 \(\langle x|y\rangle = 0\) 与 \(\langle y|x\rangle = 0\) 等价。由此推出 \(\langle y|x\rangle = \langle x|y\rangle\) 对 E 的每一对点 \(x,y\) 成立（考虑一个数 \(\lambda \in \mathbf{R}\)，满足 \(\langle \lambda x + y|x\rangle = 0\)）。
\item
  证明 \(\langle x|x\rangle\) 在所有 \(x\neq 0\) 组成的集上符号相同；必要时把 \(u\) 换成 \(-u\)，则可假设 \(\langle x|y\rangle\) 是 \(\mathbf{E}\times \mathbf{E}\) 上的正的非退化对称双线性型。d) 设 \(\mathcal{T}_0\) 是 E 的初始拓扑。证明 E 的拓扑 \(\mathcal{T}\)（由范数 \(\langle x|x\rangle^{1 / 2}\) 定义）细于拓扑 \(\mathcal{T}_0\)（注意 E 关于 \(\mathcal{T}\) 的对偶包含 E 关于 \(\mathcal{T}_0\) 的对偶 E'）。
\item
  证明：u 是从 E 到其对偶 \(E'\) 上的连续映射（对拓扑 \(\sigma(E, E')\) 与 \(\sigma(E', E)\)）。由此推出：若 E 关于初始拓扑 \(T_{0}\) 完备，则 u 关于 \(T_{0}\) 以及强拓扑 \(\beta(E', E)\) 连续（注意 u 把每个关于 \(\sigma(E, E')\) 有界的集变成关于 \(\sigma(E', E)\) 有界的集）。再推出此时拓扑 T 与 \(T_{0}\) 恒同，且 \(\omega(\mathbf{M})\) 是 \(\mathbf{M}\)（关于在 \(\mathbf{E}\) 上由型 \(\langle x|y\rangle\) 定义的希尔伯特空间结构）的正交补（cf.~I, p.~17, th. 1）。
\item
  证明：在空间 \(\ell^{1}(\mathbf{N})\)（赋以由 \(\ell^{\infty}(\mathbf{N})\) 诱导的范数）中，存在一个双射映射 \(\mathbf{M}\mapsto\omega(\mathbf{M})\)（从 \(\mathfrak{M}\) 到其自身），它具有上述性质（IV, p.~47, 习题 1）。
\end{enumerate}

¶ 19) 设 E 是无穷维复赋范空间。假设存在一个双射映射 \(\omega\)（从 E 的闭向量子空间所成之集 \(\mathfrak{M}\) 到其自身），它具有习题 18 中所列的性质。

\begin{enumerate}
\def\labelenumi{\alph{enumi})}
\item
  证明：存在一个半线性映射 \(u\)（从 E 到其对偶 \(\mathbf{E}'\) 上，关于 \(\xi \mapsto \overline{\xi}\)（\(\mathbf{C}\) 的自同构）），它在相差一个标量因子的意义下确定，且使得 \(u(\mathbf{M}) = (\omega(\mathbf{M}))^\circ\)（对每个 \(\mathbf{M} \in \mathfrak{M}\)）。（按习题 18 的方式进行；利用 IV, p.~65, 习题 16，证明 \(u\) 是关于 \(\xi \mapsto \overline{\xi}\) 的恒等自同构的半线性映射；最后证明第一种情形不会出现，因为 \(\langle x, u(x) \rangle \neq 0\)（对 \(x \neq 0\)）。）
\item
  若置 \(\langle y|x\rangle = \langle x,u(y)\rangle\)，证明 \(\langle x|y\rangle = \overline{\langle y|x\rangle}\)，且 \(\langle x|x\rangle\) 在所有 \(x\neq 0\) 组成的集上符号相同（方法同习题 18）。
\item
  最后证明：由范数 \(\langle x|x\rangle^{1 / 2}\) 定义的拓扑细于 E 上的初始拓扑 \(\mathcal{T}_0\)，且当 E 关于 \(\mathcal{T}_0\) 完备时这两个拓扑恒同；在后一情形，\(\omega (\mathbf{M})\) 是 M 在希尔伯特空间 E 中的正交补。
\end{enumerate}

\begin{enumerate}
\def\labelenumi{\arabic{enumi})}
\setcounter{enumi}{19}
\tightlist
\item
  设 \(\mathbf{E}\) 是实有限维向量空间，\(\phi\) 是从 \(\mathbf{E}\) 到其对偶 \(\mathbf{E}^*\) 上的双射线性映射。设 \(\mathbf{A}\) 是 \(\mathbf{E}\) 中的对称紧凸集，以 0 为内点；假设对每个 \(x\)（在 \(\mathbf{A}\) 的边界上），方程为 \(\langle y - x, \phi(x) \rangle = 0\) 的超平面是 \(\mathbf{A}\) 的支撑超平面。
\end{enumerate}

\begin{enumerate}
\def\labelenumi{\alph{enumi})}
\item
  设 \(f(x) = |\langle x, \phi(x) \rangle|\)，并设 \(a\) 是 \(A\) 的一个边界点，\(f(x)\) 在该点处达到其极小。证明：对每个点 \(b\)（满足 \(\langle b, \phi(a) \rangle = 0\)），我们也有 \(\langle a, \phi(b) \rangle = 0\)。（注意 \(\langle x, \phi(x) \rangle \neq 0\)（对 \(x \neq 0\)），故可设 \(f(x) = \langle x, \phi(x) \rangle \geqslant 0\)；利用如下事实：\(A\) 在点 \(a\) 处的每个支撑超平面也是由 \(f(x) \leqslant f(a)\) 所定义的集的支撑超平面。）
\item
  证明 \((x, y) \mapsto \langle x, \phi(y) \rangle\) 是对称双线性型，且 A 与由所有点 \(x\)（满足 \(f(x) \leqslant \gamma\)，其中 \(\gamma\) 为适当常数）所组成的集恒同。（对 E 的维数用归纳法论证；用 \(a\) ) 的记号，考虑方程为 \(\langle x, \phi(a) \rangle = 0\) 的超平面。）
\end{enumerate}

\begin{enumerate}
\def\labelenumi{\arabic{enumi})}
\setcounter{enumi}{20}
\tightlist
\item
  设 E 是有限维复向量空间，并设 \(\phi\) 是一个双射半线性映射（关于 C 的自同构 \(\xi \mapsto \overline{\xi}\)），从 E 到其对偶 \(E^{*}\) 上。设 \(\| x\|\) 是 E 上的一个范数，使得对所有 \(x \in E\) 有 \(|\langle x, \phi(x) \rangle| = \| x\| \cdot \| \phi(x)\|\)。证明 \((y, x) \mapsto \langle x, \phi(y) \rangle\) 在相差一个常数因子的意义下是正的非退化埃尔米特型，且 \(\langle x, \phi(x) \rangle = \gamma \| x\|^2\)（\(\gamma\) 为常数）。（按习题 20 的方式论证。）
\end{enumerate}

¶ 22) 设 E 是维数 \(\geqslant 3\) 的实赋范空间，使得对 E 中每个齐次平面 P，都存在一个从 E 到 P 上的范数为 1 的连续投影算子。证明 E 上的范数是准希尔伯特的。借助 V, p.~60 的习题 2，化归为 E 是 3 维的情形，并依次建立下列命题。

\begin{enumerate}
\def\labelenumi{\alph{enumi})}
\item
  对 E 中每个齐次平面 P，存在唯一的从 E 到 P 上的范数为 1 的连续投影算子，且此投影的核是一条齐次直线 D(P)，使得 P ↦ D(P) 是从 E 的齐次平面空间到齐次直线空间上的连续双射（GT, VI, § 3, No.~5）。
\item
  球面 \(\mathrm{D}:\| x\| = 1\)（在 \(\mathrm{E}\) 中）的每个点都是球 \(\mathbf{B}\)（\(\mathrm{E}\) 的，由 \(\| x\| \leqslant 1\) 定义）中的极点。（首先证明：若 \(x\in S\) 不是极点，则其截口 \(F_{x}\)（在 \(\mathbf{B}\) 中，II, p.~87, 习题 3）将是 2 维的，方法是考虑所有的齐次平面 \(\mathbf{P}\)（过 \(x\)）；其次，在 \(F_{x}\) 中的一点处（该点处 \(F_{x}\) 在由 \(F_{x}\) 生成的平面中只有一条支撑直线）类似地进行，以证明这一假设是矛盾的；这样的点的存在性可利用 II, p.~88, 习题 7 与 p.~88, 习题 8 建立。）
\item
  球面 \(\mathbf{S}'\)（方程为 \(\| x'\| = 1\)，在对偶 \(\mathrm{E}'\)（\(\mathrm{E}\) 的）中）的每个点都是球 \(\mathbf{B}'\)（\(\mathrm{E}'\) 的，由 \(\| x'\| \leqslant 1\) 定义）中的极点。（首先注意，对每条齐次直线 \(\mathbf{D}'\)（属于 \(\mathrm{E}'\)），存在唯一的齐次平面 \(\mathrm{P}'(\mathbf{D}')\)（在 \(\mathrm{E}'\) 中），使得对 \(\mathbf{S}' \cap \mathrm{P}'(\mathbf{D}')\) 中的每个点，\(\mathbf{B}'\) 在该点处的支撑平面（由 \(a\) 知其唯一）平行于 \(\mathbf{D}'\)；而且，映射 \(\mathbf{D}' \mapsto \mathbf{P}'(\mathbf{D}')\) 是连续的。由此推出：若 \(x' \in \mathbf{S}'\) 不是 \(\mathbf{B}'\) 中的一个极点，则其截口 \(\mathbf{F}_{x'}\)（在 \(\mathbf{B}'\) 中）的维数至少是 2；为此考虑所有的齐次直线 \(\mathbf{D}'\)（平行于 \(\mathbf{B}'\) 在点 \(x'\) 处的支撑平面）。其次，考虑一个严格凸点 \(y'\)（属于 \(\mathbf{F}_{x'}\)，II, p.~88, 习题 8），以及唯一的齐次直线 \(\mathbf{D}_0'\)（平行于 \(\mathbf{F}_{x'}\) 在点 \(y'\) 处（在由 \(\mathbf{F}_{x'}\) 生成的平面中）的支撑直线），并证明函数 \(\mathbf{D}' \mapsto \mathbf{P}'(\mathbf{D}')\) 关于 \(\mathbf{D}' = \mathbf{D}_0'\) 将不连续。）\(d)\) 证明：若 E 中三个齐次平面 \(\mathbf{P}_1, \mathbf{P}_2, \mathbf{P}_3\) 含有同一条直线 \(\Delta\)，则三条直线 \(\mathbf{D}(\mathbf{P}_1), \mathbf{D}(\mathbf{P}_2), \mathbf{D}(\mathbf{P}_3)\) 在同一齐次平面 \(\pi(\Delta)\) 中（考虑 B 在 \(\Delta\) 与 S 的一个交点处的唯一支撑平面）。应用射影几何的基本定理（A, II, § 9, 习题 16），推出存在一个双射线性映射 \(\phi\)（从 \(\mathbf{E}'\) 到 E 上），使得对所有 \(x' \in \mathbf{E}'\)，点 \(\phi(x')\) 属于直线 \(\mathbf{D}(\mathbf{P})\)，其中 \(\mathbf{P}\) 是方程为 \(\langle y, x' \rangle = 0\) 的平面。证明对每个点 \(x' \in \mathbf{S}'\)，方程为 \(\langle \phi(x'), y' - x' \rangle = 0\) 的平面是 \(\mathbf{B}'\) 在该点处的支撑平面，最后应用 V, p.~65, 习题 20 得出结论。
\end{enumerate}

\begin{enumerate}
\def\labelenumi{¶ \arabic{enumi})}
\setcounter{enumi}{22}
\item
  设 E 是维数 \(\geqslant 3\) 的复赋范向量空间，使得对 E 中每个（复）齐次平面 P，存在从 E 到 P 上的范数为 1 的连续投影算子。证明 E 上的范数是准希尔伯特的。利用 V, p.~60, 习题 2，化归为 E 在 C 上是 3 维的情形，并按习题 22 的方式进行。（对证明的 b) 部分，对每个 \(x' \in E'\)（满足 \(\|x'\| = 1\)），考虑凸集 \(G_{x'}\)（由所有 \(x \in S\) 中满足 \(\langle x, x' \rangle = 1\) 者组成），证明若 \(G_{x'}\) 不恰好是 0，则它在 R 上维数至少是 3；然后在由 \(G_{x'}\) 生成的实仿射线性流形中，考虑 \(G_{x'}\) 的一个边界点，在该点处 \(G_{x'}\) 只有唯一的（实）支撑超平面。类似地，对证明的 c) 部分，对所有 \(x \in S\)，考虑集 \(G_{x}'\)（由所有 \(x' \in S'\) 中满足 \(\langle x, x' \rangle = 1\) 者组成），并证明 \(G_{x}'\) 缩为一点；为此证明：若不然，由 \(G_{x}'\) 生成的实仿射线性流形在 R 上维数至少是 2，且含有两个在 C 上线性无关的向量。最后利用 V, p.~65, 习题 21 得出结论。）
\item
  在实赋范空间 \(\mathbf{E}\)（维数为 \(\geqslant 3\)）中，我们称向量 \(y\) 拟垂直于向量 \(x\)，如果对每个标量 \(\lambda\)，我们有 \(\| x + \lambda y \| \geqslant \| x \|\)。
\end{enumerate}

\begin{enumerate}
\def\labelenumi{\alph{enumi})}
\item
  证明：若关系「y 拟垂直于 x」关于 x、y 是对称的，则 E 上的范数是准希尔伯特的（证明 V, p.~65, 习题 22 的条件满足）。
\item
  证明：若对 E 中每个闭齐次超平面 H，都存在一个向量 \(\neq 0\) 拟垂直于 H 的所有向量，则同样的结论成立。（同一方法，把 E, III, § 2, No.~4 的 th. 2 应用于从含 P 的向量子空间到齐次平面 P 上的范数为 1 的连续投影算子，这些投影按延拓关系是线性有序的。）
\item
  证明：若对 E 中每个向量 \(x \neq 0\)，都存在一个闭超平面 H，使得 x 拟垂直于 H 的所有向量，则同样的结论成立。（化归为 E 是 3 维的情形，并对 E 的对偶应用 V, p.~65, 习题 22。）
\item
  证明：若当 z 拟垂直于 x 和 y 时，z 也拟垂直于 \(x + y\)，则同样的结论成立（应用 V, p.~65, 习题 22）。
\end{enumerate}

\begin{enumerate}
\def\labelenumi{\arabic{enumi})}
\setcounter{enumi}{24}
\tightlist
\item
  \begin{enumerate}
  \def\labelenumii{\alph{enumii})}
  \tightlist
  \item
    设 E 是实赋范空间，\(x' \neq 0\) 是 E 的对偶 E' 中的一个向量。证明：为使每个向量 \(y\)（在超平面 \(x'^{-1}(0)\) 中）都拟垂直于 \(x\)（习题 24），必须且只需 \(\langle x, x' \rangle = \| x \| \cdot \| x' \|\)。
  \end{enumerate}
\end{enumerate}

\begin{enumerate}
\def\labelenumi{\alph{enumi})}
\setcounter{enumi}{1}
\item
  由 \(a)\) 推出：对所有 \(x \neq 0\)（在 \(\mathbf{E}\) 中），存在一个闭齐次超平面 \(\mathbf{H}\)（\(\mathbf{E}\) 的），使得每个向量 \(y \in \mathbf{H}\) 都拟垂直于 \(x\)。
\item
  若 \(x, y\) 是 \(E\) 中的两个点且 \(x \neq 0\)，则存在一个标量 \(\alpha\)，使得 \(\alpha x + y\) 拟垂直于 \(x\)。
\end{enumerate}

\begin{enumerate}
\def\labelenumi{\arabic{enumi})}
\setcounter{enumi}{25}
\tightlist
\item
  若实赋范空间 E 的单位球面的所有点都是单位球的光滑点（II, p.~87, 习题 6），则称 E 是光滑的。要使这一点成立，必须且只需存在唯一的正齐次映射 f，从 E - \{0\} 到 E' - \{0\} 中，使得 \(\|f(x)\| = 1\)（对 \(\|x\| = 1\)），且 \(\langle x, f(x) \rangle = \|x\| \cdot \|f(x)\|\)。证明下列性质等价：
\end{enumerate}

α) E 是光滑的。

\(\beta)\) 对 E 中所有 \(x \neq 0\) 与所有 \(y \in \mathbf{E}\)，存在唯一的标量 \(\alpha\)，使得 \(\alpha x + y\) 拟垂直于 \(x\)。

\(\gamma)\) 对每个 \(x \in E\)，若 \(y\) 与 \(z\) 都拟垂直于 \(x\)，则 \(y + z\) 拟垂直于 \(x\)。

（为看出 \(\gamma\) ) 蕴含 \(\beta\) )，注意若 \(\alpha x + y\) 与 \(\beta x + y\) 拟垂直于 \(x\)，则 \((\alpha - \beta)x\) 拟垂直于 \(x\)。）

\begin{enumerate}
\def\labelenumi{\arabic{enumi})}
\setcounter{enumi}{26}
\item
  若实赋范空间 E 的单位球面的所有点都是单位球的严格凸点（II, p.~87, 习题 6），则称 E 是严格凸的。证明：为使 E 严格凸，必须且只需对所有 \(x \neq 0\)（在 E 中）与所有 \(y \in E\)，存在唯一的标量 \(\alpha\)，使得 \(x\) 拟垂直于 \(x + y\)。（注意映射 \(t \mapsto \| tx + y \|\) 在 R 上是凸的。）
\item
  设 \(\mathbf{E}\) 是赋范空间，\(\mathbf{E}'\) 是它的对偶。
\end{enumerate}

\begin{enumerate}
\def\labelenumi{\alph{enumi})}
\tightlist
\item
  证明：若 \(\mathrm{E}'\) 是光滑的（V, p.~66, 习题 26），则 E 是严格凸的（若 \(x\) 与 \(y\) 满足 \(x \neq y\)、\(\| x \| = \| y \| = \| \frac{1}{2}(x + y) \| = 1\)，考虑一个 \(x' \in \mathrm{E}'\)（满足 \(\| x' \| = 1\)）以及
\end{enumerate}

\[
\left\langle \frac {1}{2} (x + y), x ^ {\prime} \right\rangle = 1).
\]

\begin{enumerate}
\def\labelenumi{\alph{enumi})}
\setcounter{enumi}{1}
\tightlist
\item
  证明：若 \(\mathbf{E}'\) 是严格凸的，则 \(\mathbf{E}\) 是光滑的。
\end{enumerate}

¶ 29) 设 E 是赋范空间，E' 是它的对偶。称从 E - \{0\} 到 E' - \{0\} 中的映射 f 为支撑映射，如果它是正齐次的，且对每个 x ∈ E（满足 \textbar\textbar x\textbar\textbar{} = 1），我们有 \textbar\textbar f(x)\textbar\textbar{} = 1 与 ⟨x, f(x)⟩ = 1。为使 E 光滑（V, p.~66, 习题 26），必须且只需存在唯一的从 E - \{0\} 到 E' - \{0\} 中的支撑映射。

设 \(S\) 是 \(E\) 中的单位球面，\(S'\) 是 \(E'\) 中的单位球面，并设 \(x_0 \in S\)。下列条件等价：

α) \(x_{0}\) 是 E 中单位球的一个光滑点，

β) 存在一个支撑映射 \(f\)，它在 \(S\) 上的限制在点 \(x_0\) 处连续（当 \(S\) 赋以范数拓扑，且 \(S'\) 赋以弱拓扑 \(\sigma(E', E)\) 时）。

\(\gamma)\) 对每个 \(y \in \mathbf{E}\)，映射 \(t \mapsto \| x_0 + ty\|\) 在点 \(t = 0\) 处有导数。

（为看出 \(\alpha\) ) 蕴含 \(\beta\) )，用 E' 中单位球的弱紧性以反证法论证。为看出 \(\beta\) ) 蕴含 \(\gamma\) )，化归为 E 是 2 维的情形，并利用 \(t \mapsto \|x_0 + ty\|\) 是凸的这一事实。）

此时每个支撑映射都在点 \(x_0\) 处连续。

\begin{enumerate}
\def\labelenumi{\arabic{enumi})}
\setcounter{enumi}{29}
\tightlist
\item
  设 E 是 Banach 空间，\(E'\) 是它的强对偶，\(E''\) 是 \(E'\) 的强对偶，\(E'''\) 是 \(E''\) 的强对偶，\(E^{IV}\) 是 \(E'''\) 的强对偶。
\end{enumerate}

\begin{enumerate}
\def\labelenumi{\alph{enumi})}
\item
  假设 E 是非自反的；则存在 \(x' \in E'\) 使得 \(\| x' \| = 1\)，但关系 \(\langle x, x' \rangle = 1\) 对任何 \(x \in E\)（满足 \(\| x \| = 1\)）都不成立（IV, p.~57, 习题 25）。另一方面，存在一个点的序列 \((x_n')\)（由 \(E'\) 的点组成），使得 \(\| x_n' \| = 1\) 且强收敛于 \(x'\)，并存在一个由 E 的点组成的序列 \((x_n)\)，使得 \(\| x_n \| = 1\) 与 \(\langle x_n, x_n' \rangle = 1\)（对所有 \(n\)）（II, p.~77, 习题 4）。证明：在 \(E''\) 中，序列 \((x_n)\) 对拓扑 \(\sigma(E'', E''')\) 不收敛于任何点（注意否则它将收敛于一点 \(x \in E\)，对它有 \(\langle x, x' \rangle = 1\)）。
\item
  证明：不可能 \(x'\) 与 \(x_n'\) 都是 \(\mathrm{E}'''\) 中单位球面的光滑点（注意 \(x_n\)（看作 \(\mathrm{E}^{\mathrm{IV}}\) 的元素时）将是唯一的元素 \(x_n^{\mathrm{IV}} \in \mathrm{E}^{\mathrm{IV}}\)，它满足 \(\| x_n^{\mathrm{IV}} \| = 1\) 与 \(\langle x_n', x_n^{\mathrm{IV}} \rangle = 1\)，然后利用习题 29）。
\item
  由此得出：若 \(\mathrm{E}'''\) 是光滑的，或若 \(\mathrm{E}^{\mathrm{IV}}\) 是严格凸的，则 \(\mathrm{E}\) 必定是自反的。
\end{enumerate}

\begin{enumerate}
\def\labelenumi{\arabic{enumi})}
\setcounter{enumi}{30}
\tightlist
\item
  赋范空间 E（实的或复的）称为一致凸的，如果对每个 \(\varepsilon\)（满足 \(0 < \varepsilon < 2\)），存在 \(\delta > 0\)，使得 E 中的关系 \(\|x\| \leqslant 1\)、\(\|y\| \leqslant 1\)、\(\|x - y\| \geqslant \varepsilon\) 蕴含 \(\left\|\frac{1}{2}(x + y)\right\| \leqslant 1 - \delta\)。一致凸空间是严格凸的（V, p.~67, 习题 27）。我们称 E 是一致光滑的，如果对每个 \(\varepsilon > 0\)，存在 \(\eta > 0\)，使得关系 \(\|x\| \geqslant 1\)、\(\|y\| \geqslant 1\)、\(\|x - y\| \leqslant \eta\) 蕴含不等式 \(\|x + y\| \geqslant \|x\| + \|y\| - \varepsilon \|x - y\|\)。这等价于：对每个 \(\varepsilon > 0\)，存在 \(\rho > 0\)，使得关系 \(\|x\| = 1\)、\(\|y\| \leqslant \rho\) 蕴含不等式
\end{enumerate}

\[
\| x + y \| + \| x - y \| \leqslant 2 + \varepsilon \| y \|.
\]

一致光滑空间是光滑的（V, p.~66, 习题 26）。

\begin{enumerate}
\def\labelenumi{\alph{enumi})}
\item
  证明：若 E 一致凸，则它的强对偶 \(E'\) 一致光滑；若 E 一致光滑，则 \(E'\) 一致凸；唯一的支撑映射（V, p.~67, 习题 29）在 E 的单位球面 S 上的限制是一个从 S 到单位球面 \(S'\)（\(E'\) 的）中的映射，它关于 E 与 \(E'\) 的范数拓扑是连续的。
\item
  证明：若 E 一致凸，且 E 上的一个滤子 \(\mathfrak{F}\) 收敛于 \(x_{0}\)（对拓扑 \(\sigma(\mathrm{E}, \mathrm{E}')\)）并满足 \(\lim_{\mathfrak{F}} \|x\| = \|x_{0}\|\)，则 \(\mathfrak{F}\) 对 E 的初始拓扑收敛于 \(x_{0}\)。
\item
  证明：一致凸或一致光滑的 Banach 空间是自反的（利用 b) 与 c)，并利用 IV, p.~60, 习题 12）。（cf.~V, p.~72, 习题 14。）
\item
  把 V, p.~10 的 th. 1 的第一部分以及 V, p.~11 的 cor. 1 与 2 推广到一致凸 Banach 空间。
\end{enumerate}

\begin{enumerate}
\def\labelenumi{\arabic{enumi})}
\setcounter{enumi}{31}
\tightlist
\item
  设 \(\mathrm{E}\) 是维数 \(\geqslant 2\) 的赋范空间（实的或复的），使得对每个 \(\varepsilon\)（满足 \(0 < \varepsilon < 2\)），关系 \(\| x \| = 1\)、\(\| y \| = 1\)、\(\| x - y \| \geqslant \varepsilon\)（在 \(\mathrm{E}\) 中）蕴含不等式 \(\| \frac{1}{2}(x + y) \| \leqslant \left(1 - \frac{\varepsilon^2}{4}\right)^{1/2}\)。证明 \(\mathrm{E}\) 上的范数是准希尔伯特的。（化归为 E 是实且 2 维的情形，并按 V, p.~61 的习题 4, a) 的方式论证。）
\end{enumerate}

¶ 33) 设 E 是一致凸 Banach 空间（V, p.~67, 习题 31）。则存在一个数 θ，使得 \(\frac{3}{4} \leqslant \theta < 1\)，且使得 E 中的关系 \(\|x - y\| \geqslant \frac{1}{2} \sup (\|x\|, \|y\|)\) 蕴含 \(\left\|\frac{1}{2}(x + y)\right\| \leqslant \theta \sup (\|x\|, \|y\|)\)。

\begin{enumerate}
\def\labelenumi{\alph{enumi})}
\tightlist
\item
  设 \((x_{n})\) 是 E 中点的序列，满足 \(\| x_{n}\| \leqslant \mathbf{M}\)，且该序列对 \(\sigma (\mathrm{E},\mathrm{E}^{\prime})\) 收敛于 0。证明：若对一个指标 \(p\) 有 \(\| x_p\| \geqslant \frac{1}{2}\mathbf{M}\)，则存在 \(q > p\) 使得 \(\| x_p - x_q\| >\frac{1}{2}\mathbf{M}\)，从而 \(\left\| \frac{1}{2} (x_p + x_q)\right\| \leqslant \theta \mathbf{M}\)（用反证法论证，注意对所有 \(x^{\prime}\in \mathrm{E}^{\prime}\)（满足 \(\| x^{\prime}\| = 1\)），我们有 \(\langle x_p,x'\rangle = \lim_{n\to \infty}\langle x_p - x_n,x'\rangle\)）。
\end{enumerate}

由此推出：存在一个严格递增映射 \(\ell\)（从 \(\mathbf{N}\) 到其自身），使得 \(\left\| \frac{1}{2} (x_{\ell(2n)} + x_{\ell(2n+1)})\right\| \leqslant \mathrm{M}\theta\)，且使得若 \(x_{n}^{(1)} = \frac{1}{2} (x_{\ell(2n)} + x_{\ell(2n+1)})\)，则序列 \((x_{n}^{(1)})\) 对 \(\sigma(\mathrm{E},\mathrm{E}')\) 收敛于 0，并且 \(\| x_n^{(1)}\| \leqslant \mathrm{M}\theta\)（对所有 \(n\)）。\\
b) 证明：存在一个序列 \((x_{n_k})\)（从 \((x_{n})\) 中抽取），使得若置 \(y(k) = x_{n_k}\)，我们有下列性质：对每个整数 \(p > 1\)、每个整数 \(q < p\) 与每个整数 \(i\)（满足 \(1 \leqslant i \leqslant 2^{p - q}\)），

\[
\left\| y ((i - 1) 2 ^ {q} + 1) + y ((i - 1) 2 ^ {q} + 2) + \dots + y (i 2 ^ {q}) \right\| \leqslant \mathrm{M} \theta^ {q}.
\]

（迭代 \(a\) ) 的步骤，构造一个序列 \((x_{n}^{(k + 1)})\)（从序列 \((x_{n}^{(k)})\) 出发），其方式与 \((x_{n}^{(1)})\) 由 \((x_{n})\) 构造的方式相同；然后使用一个适当的「对角线方法」。）\(c)\) 设 \(r\) 与 \(q\) 是两个整数 \(>1\)。由 \(b\) ) 推出：若 \(r2^{q}\leqslant k\leqslant (r + 1)2^{q}\)，则

\[
\left\| x _ {n _ {1}} + x _ {n _ {2}} + \dots + x _ {n _ {k}} \right\| \leqslant (2 ^ {q} - 1) \mathrm{M} + 2 ^ {q} \mathrm{M} + (r - 1) 2 ^ {q} \mathrm{M} \theta^ {q}
\]

（把左端的和分解成若干个子集，方法是让 h 从 1 变到 \(2^{q}\)，然后从 \((j - 1)\) \(2^{q} + 1\) 变到 \(j2^{q}\)（对 \(2 \leqslant j \leqslant r\)），然后从 \(r2^{q} + 1\) 变到 k）。

\begin{enumerate}
\def\labelenumi{\alph{enumi})}
\setcounter{enumi}{3}
\tightlist
\item
  证明：对 E 中每个有界序列 \((x_{n})\)，存在一个抽取的序列 \((x_{n_{k}})\)，使得平均值的序列 \((x_{n_{1}} + \cdots + x_{n_{k}})/k\) 对 E 的初始拓扑收敛（Banach-Saks-Kakutani 定理）。（利用 E 是自反的这一事实，化归为序列 \((x_{n})\) 对 \(\sigma(\mathrm{E}, \mathrm{E}')\) 收敛于 0 的情形，并对足够大的 q 与 r 利用 c)。）
\end{enumerate}

\begin{enumerate}
\def\labelenumi{\arabic{enumi})}
\setcounter{enumi}{33}
\tightlist
\item
  设 E 是 Banach 空间，K 是凸的、有界的且对 \(\sigma(\mathrm{E},\mathrm{E}')\) 闭的子集。假设对 K 中每个序列 \((x_{n})\)，存在一个抽取的序列 \((x_{n_{k}})\)，使得平均值的序列 \((x_{n_{1}}+\cdots+x_{n_{k}})/k\) 对 \(\sigma(\mathrm{E},\mathrm{E}')\) 收敛。证明：对 E 上每个连续线性型 \(x'\)，存在 K 的一个元素 x，使得 \(\langle x,x'\rangle=\sup_{y\in K} \langle y,x'\rangle\)
\end{enumerate}

（把假设应用于一个序列 \((x_{n})\)（由 \(\mathbf{K}\) 的点组成），使得 \(\langle x_n, x'\rangle\) 趋于 \(\sup_{y\in \mathbf{K}}\langle y, x'\rangle\)）。

由此推出：若 \(\mathbf{E}\) 具有习题 33 的性质 \(d\) )，则 \(\mathbf{E}\) 是自反空间（cf.~IV, p.~57, 习题 25）。

* 35) 用 E 表示维数 \(n\) 的实有限维希尔伯特空间，S 是 E 的单位球面，\(m\) 是 S 上范数为 1 且在 E 的自同构群下不变的唯一正测度。把 S 看作一个由 \(d(x, y) = \operatorname{Arc} \cos \langle x | y \rangle\) 定义距离的度量空间。对每个 \(x \in S\) 与 \(r \geqslant 0\)，用 \(\mathbf{B}(x, r)\) 表示 S 中所有满足 \(d(x, y) \leqslant r\) 的点 y 组成的集；对 S 的每个子集 A 与每个实数 \(r \geqslant 0\)，用 \(A_r\) 表示 S 中所有满足 \(d(x, A) \leqslant r\) 的点 x 组成的集。

\begin{enumerate}
\def\labelenumi{\alph{enumi})}
\item
  给定 S 的两个闭子集 A 与 B，用 \(\delta(A, B)\) 表示由所有实数 \(r \geqslant 0\)（满足 \(A \subset B_r\) 与 \(B \subset A_r\)）组成的集的下确界。证明 \(\delta\) 是 S 的所有闭子集组成的集 \(\mathcal{F}\) 上的一个距离，且 \(\mathcal{F}\) 对此距离是紧度量空间。证明映射 \(A \mapsto m(A)\)（从 \(\mathcal{F}\) 到 \(\mathbf{R}\) 中）是上半连续的。
\item
  设 \(x_0\) 是 \(S\) 中的一点，设 \(H\) 是 \(E\) 中正交于 \(x_0\) 的超平面，\(x_1\) 是 \(H\) 中的一点，\(\gamma\) 是联结 \(x_0\) 与 \(-x_0\) 且过 \(x_1\) 的圆弧，即 \(S\) 中所有形如 \(x_0 \sin \theta + x_1 \cos \theta\) 的点（\(|\theta| \leqslant \pi / 2\)）组成的集。对每个 \(y \in \gamma\)，令 \(H_y = H + y\) 与 \(S_y = S \cap H_y\)；用 \(m_y\) 表示 \(S_y\) 上范数为 1 的唯一正测度，它在 \(E\) 的所有使 \(x_0\) 不动的自同构组成的群下不变。
\end{enumerate}

设 A 是 S 的一个闭子集，\(\gamma'\) 是 \(\gamma\) 中所有使 \(A \cap S_{p}\) 非空的点组成的集。对每个 \(y \in \gamma'\)，存在唯一的实数 \(r(y)\)，使得 \(0 \leqslant r(y) \leqslant \pi\) 且 \(m_{y}(A \cap S_{y}) = m_{y}(B(y, r(y)) \cap S_{y})\)；用 \(s_{\gamma}(A)\) 表示诸集 \(B(y, r(y)) \cap S_{y}\)（当 y 取遍 \(\gamma'\) 时）的并。证明 \(s_{\gamma}(A)\) 是闭的，且 \(m(A) = m(s_{\gamma}(A))\)。

\begin{enumerate}
\def\labelenumi{\alph{enumi})}
\setcounter{enumi}{2}
\tightlist
\item
  对 S 的每个闭子集 A，\(r(\mathbf{A})\)（由所有实数 \(r \geqslant 0\)（对之存在 \(x \in S\) 满足 \(\mathbf{A} \subset \mathbf{B}(x, r)\)）组成的集的下确界）称为 A 的半径。用 \(\mathbf{M}(\mathbf{A})\) 表示 S 的所有闭子集 C（满足 \(m(\mathbf{C}) = m(\mathbf{A})\) 且 \(m(\mathbf{C}_{\varepsilon}) \leqslant m(\mathbf{A}_{\varepsilon})\)（对每个 \(\varepsilon > 0\)））组成的集。证明：对 S 的闭子集的每个对 (A, B)，下列条件等价：
\end{enumerate}

\begin{enumerate}
\def\labelenumi{(\roman{enumi})}
\item
  \(m(\mathbf{A}) = m(\mathbf{B})\) 且 \(\mathbf{B}\) 形如 \(\mathbf{B}(x,r)\)（其中 \(x\in \mathbf{S}\) 与 \(r\geqslant 0\)）
\item
  \(\mathbf{B} \in \mathbf{M}(\mathbf{A})\) 且 \(r(\mathbf{B}) \leqslant r(\mathbf{C})\) 对每个子集 \(\mathbf{C}\)（\(\mathbf{A}\) 的，属于 \(\mathbf{M}(\mathbf{A})\)）成立。（对 \(n\) 用归纳法，由 \(b\) ) 推出 \(s_{\gamma}(\mathbf{A})\) 属于 \(\mathbf{M}(\mathbf{A})\)，对每个闭子集 \(\mathbf{A}\)（\(\mathbf{S}\) 的）都如此；若 \(r > 0\) 满足 \(\mathbf{A} \subset \mathbf{B}(x_1, r)\)，证明 \(\mathbf{B}(x_1, r)\) 在 \(\mathbf{S}\) 中的边界上的每个点，凡属于 \(s_{\gamma}(\mathbf{A})\) 者，也属于 \(\mathbf{A}\)。）
\end{enumerate}

* 36) 记号与习题 35 中的相同。

\begin{enumerate}
\def\labelenumi{\alph{enumi})}
\item
  设 \(a\) 是 \(\mathbf{E}\) 中范数为 1 的向量，\(\mathrm{K}_{\varepsilon}\) 是所有 \(x \in \mathbf{S}\)（满足 \(|\langle x|a\rangle| \leqslant \sin \varepsilon\)）组成的集，\(\mathrm{L}_{\varepsilon}\) 是所有 \(x \in \mathbf{S}\)（满足 \(d(x, \mathrm{S}_a) \geqslant \varepsilon\)）组成的集（其中 \(\mathrm{S}_a\) 是 \(\mathbf{S}\) 中所有正交于 \(a\) 的点组成的集）。证明：对足够小的 \(\varepsilon > 0\)，我们有 \(m(\mathrm{K}_{\varepsilon}) \geqslant 4e^{-n\varepsilon^2 / 2}\) 与 \(m(\mathrm{L}_{\varepsilon}) \leqslant 4e^{-n\varepsilon^2 / 2}\)（注意测度 \(m\) 在映射 \(x \mapsto \langle x|a\rangle\)（从 \(\mathbf{S}\) 到区间 \([-1, 1]\)（\(\mathbf{R}\) 的）中）下的像形如 \(c_n(1 - t^2)^{(n-3)/2}dt\)，其中 \(c_n > 0\) 为适当的常数）。
\item
  设 f 是从 S 到 R 中的连续映射，\(\mathbf{M}(f)\) 是一个实数，使得所有 \(x \in S\)（对之 \(f(x) \leqslant \mathbf{M}(f)\)（相应地 \(f(x) \geqslant \mathbf{M}(f)\)）者）组成的集对 m 有测度 \(\geqslant \frac{1}{2}\)。设 B 是所有 \(x \in S\)（满足 \(f(x) = \mathbf{M}(f)\)）组成的集。由 a) 推出：对每个足够小的 \(\varepsilon > 0\)，所有 \(x \in S\)（满足 \(d(x, B) \geqslant \varepsilon\)）组成的集对 m 的测度至多等于 \(4e^{-n\varepsilon^{2}/2}\)。
\item
  对每个 \(\varepsilon > 0\)，令 \(h(n, \varepsilon)\) 为最小的整数 \(h \geqslant 1\)，对它存在 S 中的点 \(x_1, \ldots, x_h\) 使得 \(S = \bigcup_{i=1}^{h} B(x_i, \varepsilon)\)；证明 \(\lim_{\varepsilon \to 0} (\log h(n, \varepsilon)) / |\log \varepsilon| = n\)。
\item
  回顾 E 是维数 n 的实希尔伯特空间。设 k 是一个正整数，\(\varepsilon\)、\(\varepsilon'\) 是两个严格正的数，满足 \(4h(k, \varepsilon) < e^{n\varepsilon'^{2}/2}\)。设 f 是从 S 到 R 中的映射，使得对 S 中所有 x, y 有 \(|f(x) - f(y)| \leqslant \|x - y\|\)，并设 \(\mathbf{M}(f)\) 是满足 b) 中所述关系的实数。证明：存在 E 的一个维数为 k 的向量子空间 F，满足下列条件：对每个 \(x \in F \cap S\)，存在 \(F \cap S\) 中的一点 y，使得 \(\|x - y\| < \varepsilon\) 且 \(|f(y) - \mathbf{M}(f)| \leqslant \varepsilon'\)。
\end{enumerate}

* 37) 设 \(\mathrm{E}\) 是维数 \(n\) 的实希尔伯特空间，并设 \(\gamma\) 是 \(\mathrm{E}\) 上的一个正测度，使得 \(\int_{\mathrm{E}} e^{i\langle x|y\rangle} d\gamma(y) = \exp(-\pi \|x\|^2/2)\) 对所有 \(x \in \mathrm{E}\) 成立（INT, IX, § 6, No.~5）。设 \(m\) 是 E 的单位球面 S 上范数为 1 且在 E 的自同构群下不变的唯一正测度。

\begin{enumerate}
\def\labelenumi{\alph{enumi})}
\item
  设 \(p\) 是 \(\mathbf{E}\) 上的连续函数，满足 \(p(t,x) = t.p(x)\)（对所有 \(x\in S\) 及所有正实数 \(t\)）。证明 \(\int_{\mathbf{E}}pd\gamma = c_n\int_{\mathbf{S}}pdm\)，其中 \(c_{n} = \pi^{1 / 2}\Gamma (n / 2) / \Gamma ((n + 1) / 2)\)。
\item
  证明：存在一个常数 \(\mathbf{C} > 0\)（与 \(n\) 无关），使得
\end{enumerate}

\[
\int_ {E} \sup _ {1 \leqslant i \leqslant n} | \langle x | e _ {i} \rangle | d \gamma (x) \geqslant C. (\log n) ^ {1 / 2}
\]

对 E 的每个标准正交基 \((e_1, \dots, e_n)\)。

\begin{enumerate}
\def\labelenumi{\alph{enumi})}
\setcounter{enumi}{2}
\tightlist
\item
  设 \(\varepsilon > 0\)，并设 \(k\) 是一个正整数。由 \(b\) ) 与习题 12（V, p.~63）及习题 36, \(d\) ) 推出：若 \(n\) 足够大，则对每个维数为 \(n\) 的实赋范空间 V，存在 V 的一个维数为 \(k\) 的向量子空间 W，满足下列性质：存在一个实希尔伯特空间 \(\mathbf{W}_1\)（维数为 \(k\)），以及一个双射线性映射 \(u\)（从 W 到 \(\mathbf{W}_1\) 上），使得 \(\sup (\| u \|, \| u^{-1} \|) \leqslant 1 + \varepsilon\)。
\end{enumerate}

